\documentclass[11pt]{article}

\usepackage[T1]{fontenc}
\usepackage[utf8]{inputenc}
\usepackage{lmodern}
\usepackage[margin=1in]{geometry}
\usepackage{microtype}
\usepackage{amsmath,amssymb,mathtools,bm}
\usepackage{booktabs,tabularx,array,longtable}
\usepackage{enumitem}
\usepackage{xcolor}
\usepackage{listings}
\usepackage{hyperref}

\definecolor{codebg}{RGB}{247,247,247}
\definecolor{codeframe}{RGB}{205,205,205}
\definecolor{codecomment}{RGB}{0,105,75}
\definecolor{codekeyword}{RGB}{25,65,145}

\lstdefinestyle{base}{
  basicstyle=\ttfamily\footnotesize,
  columns=fullflexible,
  keepspaces=true,
  showstringspaces=false,
  breaklines=true,
  breakatwhitespace=false,
  frame=single,
  rulecolor=\color{codeframe},
  backgroundcolor=\color{codebg},
  numbers=left,
  numberstyle=\tiny\color{gray},
  numbersep=7pt,
  xleftmargin=1.5em,
  framexleftmargin=1em,
  aboveskip=0.8em,
  belowskip=0.8em
}

\lstdefinestyle{cpp}{
  style=base,
  language=C++,
  keywordstyle=\color{codekeyword}\bfseries,
  commentstyle=\color{codecomment},
  stringstyle=\color{purple!65!black}
}

\lstdefinestyle{data}{
  style=base,
  language={}
}

\hypersetup{
  colorlinks=true,
  linkcolor=blue!55!black,
  urlcolor=blue!65!black,
  citecolor=blue!55!black,
  pdftitle={Topology and Remeshing in ParaDiS},
  pdfsubject={Line-by-line implementation guide}
}

\setlist{nosep,leftmargin=*}
\setlist[description]{style=nextline,leftmargin=0pt,labelindent=0pt,
                      labelwidth=0pt,labelsep=0pt,
                      font=\normalfont\bfseries}
\setcounter{tocdepth}{2}
\allowdisplaybreaks
\emergencystretch=2em
\newcommand{\code}[1]{\texorpdfstring{\nolinkurl{#1}}{#1}}
\newcommand{\srcfile}[1]{\texorpdfstring{\nolinkurl{#1}}{#1}}
\newcommand{\vect}[1]{\bm{#1}}

\title{Topology and Remeshing in ParaDiS\\
       \large Execution Flow and Detailed Code Guide}
\author{ParaDiS-Extended Development Notes}
\date{August 2026}

\begin{document}

\maketitle

\begin{abstract}
ParaDiS evolves a graph of dislocation nodes and segments. Force integration
moves the graph geometrically; topology code changes its connectivity; and
remeshing changes its discretization while preserving the physical network.
This guide follows those operations from the outer time-step driver through
node splitting, node merging, collision handling, distributed-operation
replay, mesh coarsening, and mesh refinement. It is written against the
repository state at commit \code{cca5ff2}.

The code excerpts are included directly from the source tree and retain their
original line numbers. Each selected executable block is followed by a
line-by-line explanation. Large, repetitive material-specific routines are
described at their dispatch and mutation boundaries rather than copied five
times. The source files remain authoritative when a future change moves a
line or changes a condition.
\end{abstract}

\tableofcontents
\newpage

\section{Scope and reading convention}

\subsection{What topology and remesh mean in ParaDiS}

In this guide, \emph{topology} means an operation that changes the graph: an
arm is moved, a node is created or removed, a segment endpoint is relinked,
or a Burgers vector/glide plane is changed. These events arise from multi-node
separation, collision/annihilation, cross slip, segment splintering,
HCP-specific splitting, and, for free surfaces, surface operations.

\emph{Remeshing} is the later, dedicated rediscretization pass. It removes an
eligible degree-two node to coarsen a short or collapsing local shape, or
inserts a degree-two node to refine a long or highly curved segment. It uses
the same \code{SplitNode()} and \code{MergeNode()} primitives as topology
events, but it has its own selection criteria and MPI synchronization pass.

\subsection{How to read the listings}

Each listing shows exact source line numbers, not document-local line
numbers. One explanation item may cover a contiguous group of lines when the
lines together implement one expression, declaration, loop, or preprocessor
guard. Braces, blank lines, and comments are explained with their surrounding
code rather than restated mechanically. This gives the reader a practical
line-by-line guide without pretending that an isolated brace has independent
behavior.

\begin{longtable}{@{}p{0.40\linewidth}p{0.54\linewidth}@{}}
\toprule
File & Responsibility \\
\midrule
\endhead
\srcfile{src/ParadisStep.cc} & Places topology and remesh in each accepted outer cycle. \\
\srcfile{src/Topology.cc} & Shared mutation primitives, serial multi-node evaluation, ownership checks, duplicate-link cleanup. \\
\srcfile{src/Remesh.cc} & Generic remesh dispatcher, force estimates, and common mesh coarsening. \\
\srcfile{src/RemeshRule_2.cc} & Rule 2 refinement, which inserts at a segment midpoint. \\
\srcfile{src/RemeshRule_3.cc} & Rule 3 refinement, which may place a degree-two refinement point on a fitted arc. \\
\srcfile{src/CommSendRemesh.cc}\newline
\srcfile{src/FixRemesh.cc}\newline
\srcfile{src/RemoteOps.cc} & Pack, send, and replay graph mutations across MPI domains. \\
\srcfile{src/SplitMultiNodesParallel.cc}\newline
\srcfile{src/SMN*.cc} & Distributed force evaluation for parallel multi-node splitting. \\
\bottomrule
\end{longtable}

\subsection{One-cycle mental model}

\begin{center}
\fbox{\parbox{0.92\linewidth}{
Integrate positions and velocities \(\rightarrow\) initialise topology
exemptions \(\rightarrow\) identify/split multi-nodes \(\rightarrow\)
cross slip \(\rightarrow\) collision and special material events
\(\rightarrow\) distribute/replay those changes \(\rightarrow\)
remove duplicate links \(\rightarrow\) coarsen/refine the mesh
\(\rightarrow\) distribute/replay remesh changes \(\rightarrow\)
remove any orphan nodes.
}}
\end{center}

The order is intentional. Topological physics events run first; remesh then
cleans up the graph they leave behind. Two operation-list exchanges are normal
in a parallel simulation: one after collision/topology work, and another
after remeshing.

\section{Data model, controls, and invariants}

\subsection{Node state and topology flags}

\lstinputlisting[style=cpp,firstline=24,lastline=205,firstnumber=24]{../include/Node.h}

\paragraph{Line-by-line explanation.}
\begin{description}
\item[Lines 24--50] Constraints are bit flags. Pin bits identify fixed
coordinates, \code{SURFACE_NODE} identifies a free-surface node, and
\code{CORNER_NODE} marks a physical degree-two corner that remesh should not
normally collapse. The topology implementation currently treats any pin bit
as a fully pinned node; the warning at lines 30--35 is operational.
\item[Lines 72--120] The macros make bitfield tests readable. For example,
\code{HAS_ANY_OF_CONSTRAINTS(node->constraint, PINNED_NODE)} is the normal
test for a node that cannot be deleted or freely repositioned.
\item[Lines 123--131] These are transient per-cycle flags.
\code{NO_COLLISIONS} prevents a node touched by one event from being reused
by another collision in the same cycle. \code{NO_MESH_COARSEN} similarly
prevents overlapping coarsening. \code{NODE_RESET_FORCES} says cached
force/velocity values are not reliable enough for geometry-sensitive work.
\item[Lines 141--153] A node owns position, total force, velocity, and
history. The \code{old*} and \code{currv*} fields support integration and
plastic-strain bookkeeping, so topology must preserve or deliberately update
them.
\item[Lines 158--171] \code{myTag} is the globally meaningful identity.
\code{nbrTag} identifies adjacent nodes. Per-arm force, Burgers-vector, and
glide-plane arrays are the segment data stored at this endpoint.
\item[Lines 185--205] \code{constraint} is persistent state.
\code{cellIdx}, \code{cell2Idx}, and \code{cell2QentIdx} connect a node to
spatial acceleration structures. \code{native} distinguishes local ownership
from a ghost replica.
\end{description}

\subsection{Public topology and remesh interfaces}

\lstinputlisting[style=cpp,firstline=16,lastline=101,firstnumber=16]{../include/Topology.h}

\lstinputlisting[style=cpp,firstline=20,lastline=74,firstnumber=20]{../include/Remesh.h}

\paragraph{Line-by-line explanation.}
\begin{description}
\item[\srcfile{Topology.h}, lines 16--31] Merge statuses are bitwise. A caller
must test \code{MERGE_SUCCESS} rather than compare the full result: success
may coexist with \code{MERGE_NO_REPOSITION} or \code{MERGE_NODE_ORPHANED}.
\item[\srcfile{Topology.h}, lines 36--55] \code{SplitNode()} returns a simple
success/failure value. The operation classes identify separation, collision,
and remesh invocations so common primitives can enforce suitable rules.
\item[\srcfile{Topology.h}, lines 59--100] The header separates high-level
selection routines such as \code{SplitMultiNodes()} from low-level mutation
routines such as \code{SplitNode()} and \code{MergeNode()}. A new feature
should normally select an event at the high level and use these primitives.
\item[\srcfile{Remesh.h}, lines 20--24] A candidate event starts
\code{REMESH_PENDING}; an executed event becomes \code{REMESH_COMMITTED};
overlapping or invalid candidates become \code{REMESH_REJECTED}.
\item[\srcfile{Remesh.h}, lines 30--59] Coarsen/refine records store
\code{Tag_t}s instead of pointers because a prior event can delete or move a
node. A refinement record also stores estimated force values for the two
replacement segments. Rule 3 uses \code{haveNewPos} and \code{newPos} to
carry its fitted-arc insertion point from selection to commit.
\item[\srcfile{Remesh.h}, lines 64--72] \code{MeshCoarsen()} is shared by
rules 2 and 3. The rule functions differ principally in \code{MeshRefine()}
placement logic.
\end{description}

\subsection{Control-file parameters and derived thresholds}

\lstinputlisting[style=cpp,firstline=227,lastline=261,firstnumber=227]{../include/Param.h}

\lstinputlisting[style=cpp,firstline=579,lastline=627,firstnumber=579]{../src/Initialize.cc}

\paragraph{Line-by-line explanation.}
\begin{description}
\item[\srcfile{Param.h}, lines 227--243] \code{minSeg} and \code{maxSeg}
bound target discretization. \code{rann} is the annihilation/contact distance
and is reused by collisions, multi-node separation, and remesh safety guards.
The internal geometric thresholds are \code{remeshAreaMin} and
\code{remeshAreaMax}.
\item[\srcfile{Param.h}, lines 244--261] \code{splitMultiNodeFreq} throttles
multi-node work. \code{splitMultiNodeAlpha} supplies the noise penalty used
when comparing trial split power. \code{split3node} enables the special BCC
binary-junction case; the final controls select parallel evaluation and
uniform versus calculated junction length.
\item[\srcfile{Initialize.cc}, lines 579--602] If absent, position tolerance
becomes one quarter of core radius and \code{rann} becomes twice the
tolerance. These defaults propagate into topology distances.
\item[\srcfile{Initialize.cc}, lines 604--620] Minimum area starts as
\(2\,rTol\,maxSeg\), is tightened by the equilateral-triangle area implied by
\code{minSeg} when provided, and then produces maximum area. The formulas
make curvature criteria consistent with length and numerical tolerances.
\item[\srcfile{Initialize.cc}, lines 622--627] When \code{minSeg} is absent,
the code derives it from minimum area. A control file can therefore specify
\code{maxSeg}/\code{rTol} and obtain a consistent lower length limit.
\end{description}

\begin{lstlisting}[style=data]
# Typical controls; values must suit material, units, and cell size.
maxSeg                    = 200.0
minSeg                    = 50.0
rann                      = 2.5
remeshRule                = 2
splitMultiNodeFreq        = 1
splitMultiNodeAlpha       = 1.0e-3
split3node                = 1
useParallelSplitMultiNode = 0
useUniformJunctionLen     = 0
collisionMethod           = 2
\end{lstlisting}

\section{Where operations occur in one ParaDiS cycle}

\lstinputlisting[style=cpp,firstline=380,lastline=493,firstnumber=380]{../src/ParadisStep.cc}

\paragraph{Line-by-line explanation.}
\begin{description}
\item[Lines 380--389] Output and the just-finished integration work occur
before this excerpt. The first topology-specific action is
\code{InitTopologyExemptions()}, which clears temporary exclusions and repairs
stale corner marking.
\item[Lines 391--395] \code{ClearOpList()} begins a fresh serialized change
log. \code{SortNodesForCollision()} builds the fine \code{cell2} queues used
to limit collision searches. It must happen before collision operations can
delete nodes.
\item[Lines 396--401] With \code{FIX_PLASTIC_STRAIN}, correction storage is
reset before any operation changes swept area or a nodal position.
\item[Lines 403--404] Exactly one multi-node implementation runs. The serial
path evaluates trial splits locally; the parallel path distributes force
contributions before committing local results.
\item[Lines 406--414] Surface splitting is conditional on \code{_ARLFEM};
\code{CrossSlip()} is a material dispatcher. Both can change planes and
positions and participate in topology synchronization.
\item[Lines 416--436] Collisions are handled after multi-node/cross-slip
events. Segment splintering, HCP splitting, and optional particle collisions
then append further mutations to the same operation list.
\item[Lines 439--452] The first \code{CommSendRemesh()}/\code{FixRemesh()}
pair broadcasts physics-driven topology changes and replays neighboring
domains' changes locally.
\item[Lines 454--482] Burgers conservation is checked after replay. A final
duplicate-link pass is needed because separately valid remote changes can
jointly make two links connect the same pair of nodes.
\item[Lines 484--493] After topology is stable, \code{Remesh()} performs the
independent coarsen/refine pass and owns a second clear/send/fix sequence.
\end{description}

\subsection{Collision-search preparation}

\lstinputlisting[style=cpp,firstline=57,lastline=130,firstnumber=57]{../src/SortNodesForCollision.cc}

\lstinputlisting[style=cpp,firstline=265,lastline=425,firstnumber=265]{../src/SortNodesForCollision.cc}

\paragraph{Line-by-line explanation.}
\begin{description}
\item[Lines 57--103] The first call computes a permanent fine-grid size.
\code{maxSeparation = 1.1*maxSeg + rann} ensures candidates that could
interact lie in the same or a neighboring \code{cell2}.
\item[Lines 104--130] The code uses an integral number of fine cells per
ordinary cell and caps it at 20 in each direction. It records the global
number of body cells for later periodic-image mapping.
\item[Lines 265--280] A head array represents queues; a compact second array
represents queue entries. The latter has one potential entry per local node
key.
\item[Lines 282--343] Each node is folded to the image nearest the current
domain centre, mapped to \code{cell2}, and clamped into the native range when
needed. This keeps periodic images from creating unbounded search locations.
\item[Lines 394--425] The node's starting \code{cell2Idx} is stored, then a
critical section prepends it to its queue and stores the queue-entry index.
A deleted node can later have its queue pointer cleared without breaking
traversal of the remaining queue.
\end{description}

\section{The shared topology mutation layer}

\subsection{Ownership is the primary parallel invariant}

\lstinputlisting[style=cpp,firstline=8,lastline=34,firstnumber=8]{../src/Topology.cc}

\paragraph{Line-by-line explanation.}
\begin{description}
\item[Lines 8--15] The comments establish the reason for strict ownership:
all domains can see ghosts, but only a deterministic owner may mutate a shared
graph entity.
\item[Lines 17--23] Owning a segment permits changing its endpoint tag,
per-segment Burgers vector, force, and glide plane, and permits deletion.
\item[Lines 25--29] Owning a node permits changing its local state, but
deleting it is only allowed once all arms are gone. That prevents one domain
from silently removing a node still referenced by another.
\item[Lines 31--34] The implementation collects helpers that enforce these
rules. Callers should not manually edit neighbor arrays.
\end{description}

\lstinputlisting[style=cpp,firstline=2208,lastline=2280,firstnumber=2208]{../src/Util.cc}

\paragraph{Ownership schedule.}
\begin{description}
\item[Lines 2208--2251] \code{DomainOwnsSeg()} assigns local ownership
immediately when both endpoints share the same domain.
\item[Lines 2252--2267] For separation/collision across domains, the lower
domain owns a segment on even cycles and the higher domain owns it on odd
cycles.
\item[Lines 2268--2273] Remesh intentionally uses the inverse cross-domain
rule. Alternation distributes work and lets an ownership-blocked operation
become eligible on a later cycle instead of making one rank permanently
responsible.
\item[Lines 2275--2280] Invalid operation classes are treated as programmer
errors rather than being given an unsafe default owner.
\end{description}

\subsection{Per-cycle exemptions and temporary mobility}

\lstinputlisting[style=cpp,firstline=327,lastline=498,firstnumber=327]{../src/Topology.cc}

\paragraph{Line-by-line explanation.}
\begin{description}
\item[Lines 327--368] The routine clears \code{NO_COLLISIONS} and
\code{NO_MESH_COARSEN} on native and ghost nodes. It removes
\code{CORNER_NODE} if a node is no longer degree two, preventing stale
protection of a later unrelated configuration.
\item[Lines 390--415] \code{EvaluateMobility()} zeroes an orphan's velocity
instead of calling a mobility law that may require arms.
\item[Lines 417--428] Mobility is not evaluated until every neighbor can be
resolved locally. This protects material laws that examine neighbor data.
\item[Lines 430--476] The selected mobility law is called, then glide-plane
constraints and the optional shear-velocity cap are applied to the returned
velocity.
\item[Lines 488--498] An optional NaN/Inf check turns bad trial mobility into
a source-local fatal error rather than allowing corrupt state to propagate.
\end{description}

\subsection{\code{SplitNode()}: one primitive for separation and refinement}

\lstinputlisting[style=cpp,firstline=2465,lastline=2583,firstnumber=2465]{../src/Topology.cc}

\paragraph{Line-by-line explanation.}
\begin{description}
\item[Lines 2465--2487] The signature receives desired positions/velocities,
arm indices to move, an operation class, and output pointers for the two
post-split nodes. Outputs are first nulled so every failure is unambiguous.
\item[Lines 2489--2497] Only a native node may be split by this domain.
Moving zero arms is a fatal programming error because it would create an
invalid empty node.
\item[Lines 2499--2510] A one-arm split normally bisects a segment during
remesh. The source arm's glide plane is saved so both child segments inherit
it instead of receiving a newly guessed plane.
\item[Lines 2512--2530] Arm removal renumbers the array. The temporary arm
list is sorted in descending order, guaranteeing that deleting a later index
never invalidates an earlier index still to be processed.
\item[Lines 2532--2541] The operation is permitted only if this domain owns
each selected segment. Endpoint relinking is a segment mutation.
\item[Lines 2543--2583] If requested arms are not all owned, the routine tries
the complementary partition. \code{switchedArms} preserves logical
\code{splitNode1}/\code{splitNode2} semantics; if neither partition is
locally owned, the split fails without graph mutation.
\end{description}

\lstinputlisting[style=cpp,firstline=2585,lastline=2712,firstnumber=2585]{../src/Topology.cc}

\paragraph{Line-by-line explanation.}
\begin{description}
\item[Lines 2585--2603] \code{GetNewNativeNode()} obtains a reusable record.
Existing arm arrays are freed because this primitive requires a zero-arm
starting state. The new node becomes native, unconstrained, and has zero old
velocity.
\item[Lines 2605--2617] In a free-surface build, the optional flag duplicates
surface constraint/normal data. Ordinary bulk remesh does not use this branch.
\item[Lines 2619--2647] The routine maps the original/new records to
\code{splitNode1}/\code{splitNode2} according to whether it complemented the
selected arm set. It writes caller-chosen positions and velocities and copies
old position to the new node.
\item[Lines 2653--2684] A pinned original remains on the \code{pos1} side;
surface properties similarly remain there unless explicitly duplicated. This
prevents a generic split from accidentally freeing an anchored endpoint.
\item[Lines 2686--2712] A global operation serializes a
\code{REM_OP_SPLIT_NODE} record plus a coordinate reset for the original
node. The reset is needed because a caller may have moved that node.
\end{description}

\lstinputlisting[style=cpp,firstline=2714,lastline=2860,firstnumber=2714]{../src/Topology.cc}

\paragraph{Line-by-line explanation.}
\begin{description}
\item[Lines 2714--2746] For every selected arm, the code gets the neighbor,
saves endpoint force, subtracts old contribution, relinks the neighbor to the
new node, inserts the matching new arm, removes the old arm, and restores
force. Helper calls keep both endpoints and the operation log consistent.
\item[Lines 2748--2756] The temporary arm list is released and the new node
is assigned to a spatial cell before later force work can find it.
\item[Lines 2758--2778] The routine sums Burgers vectors attached to the new
side. Unless both nodes are surface nodes, a nonzero sum requires a connecting
segment for Burgers conservation.
\item[Lines 2780--2799] The long comment explains a pinned-node corner case.
The nonzero Burgers-sum test is the usual condition for adding the link.
\item[Lines 2801--2848] A refinement split copies the original plane. General
separation computes a plane from line direction and Burgers vector, uses
velocity as a fallback, then chooses an allowed screw plane if undefined.
\item[Lines 2850--2860] Opposite Burgers vectors are inserted at the two
endpoints of the connecting segment. The routine leaves its force to the
caller, which knows whether a saved or estimated force is appropriate.
\end{description}

\subsection{\code{MergeNode()}: common coarsen and collision primitive}

\lstinputlisting[style=cpp,firstline=2879,lastline=2929,firstnumber=2879]{../src/Topology.cc}

\lstinputlisting[style=cpp,firstline=2976,lastline=3117,firstnumber=2976]{../src/Topology.cc}

\paragraph{Line-by-line explanation.}
\begin{description}
\item[Lines 2879--2929] \code{MergeNodePreCondition()} is a fast version of
the deletion eligibility test. It returns which supplied endpoint, if either,
is a legal node to delete.
\item[Lines 2976--2995] \code{MergeNode()} starts pessimistically: it returns
no merged node and \code{MERGE_NOT_PERMITTED} until a legal deletion target
is found.
\item[Lines 2996--3024] \code{node1Deletable} requires local ownership, no
pin, compatible surface condition, no per-cycle collision/coarsen exemption,
and local ownership of every incident segment.
\item[Lines 3027--3057] If node 1 cannot be deleted, the same test is applied
to node 2. At most one legal choice is needed; if neither is deletable,
nothing is changed.
\item[Lines 3059--3117] A prospective double link is checked before mutation.
If cleanup of that double link would be owned remotely, the merge is rejected.
Never create a graph state that the current domain cannot repair.
\end{description}

\lstinputlisting[style=cpp,firstline=3119,lastline=3268,firstnumber=3119]{../src/Topology.cc}

\paragraph{Line-by-line explanation.}
\begin{description}
\item[Lines 3119--3135] The chosen survivor is returned and marked
successful. If it is local and position-compatible, it is moved to the
requested merge position.
\item[Lines 3136--3151] With plastic-strain correction enabled, collision or
remesh merging accounts for swept area from both nodes before coordinates
change.
\item[Lines 3154--3186] The target coordinate is written and, for global
events, logged for remote replay. If constrained or remote, the merge still
proceeds but reports \code{MERGE_NO_REPOSITION}.
\item[Lines 3188--3229] Any direct target/dead link is deleted first. Each
remaining dead-node arm is transferred: the neighbor is relinked, the dead
endpoint arm is deleted, the survivor receives the same Burgers vector/plane,
and its force is restored.
\item[Lines 3231--3241] All dead-node constraints are inherited by the
survivor, then the now-arm-free dead node is removed.
\item[Lines 3243--3268] Duplicate links are combined or annihilated. If that
leaves an orphan, a local orphan is deleted immediately and a remote orphan
is reported for later cleanup.
\end{description}

\subsection{Duplicate links and orphan nodes}

\lstinputlisting[style=cpp,firstline=91,lastline=109,firstnumber=91]{../src/Topology.cc}

\lstinputlisting[style=cpp,firstline=523,lastline=708,firstnumber=523]{../src/Topology.cc}

\paragraph{Line-by-line explanation.}
\begin{description}
\item[Lines 91--109] \code{RemoveOrphanedNodes()} scans local node keys and
deletes a node with zero arms. It deliberately does not message other
domains; the subsequent ghost exchange publishes the corrected list.
\item[Lines 523--572] \code{RemoveDoubleLinks()} finds repeated neighbor tags
on one local node. It records relevant arm data before deleting or combining
them.
\item[Lines 573--593] If the Burgers vectors cancel, both duplicate arms are
removed. This is topological annihilation of redundant opposite segments.
\item[Lines 603--695] If they do not cancel, the vectors are combined and a
valid glide plane is retained or recalculated. The routine updates forces and
the remote-operation stream as requested.
\item[Lines 696--708] A resulting local orphan is deleted. The return status
lets callers record that an operation made an orphan, even when it is remote.
\end{description}

\section{Topology-event selection}

\subsection{Serial multi-node splitting}

\lstinputlisting[style=cpp,firstline=1354,lastline=1510,firstnumber=1354]{../src/Topology.cc}

\paragraph{Line-by-line explanation.}
\begin{description}
\item[Lines 1354--1382] The routine declares candidate partitions, backups,
force records, mobility arguments, and the best power-release candidate.
\item[Lines 1385--1400] The frequency guard lets diagnostics remain active
while skipping costly tests on other cycles. Initial desired separation is
\(2rann+\epsilon\); \code{vNoise} penalizes candidates whose apparent
advantage is within time-step/mobility noise.
\item[Lines 1424--1461] Native nodes with at least four arms are candidates.
A three-arm node is admitted only when \code{split3node} is enabled for BCC
and \code{BCC_binary_junction_node()} identifies a non-planar binary junction.
\item[Lines 1463--1510] Surface and pinned nodes are excluded. A short arm
also excludes the candidate because splitting near an already tiny segment can
create high-velocity numerical debris.
\end{description}

\lstinputlisting[style=cpp,firstline=1512,lastline=1642,firstnumber=1512]{../src/Topology.cc}

\paragraph{Line-by-line explanation.}
\begin{description}
\item[Lines 1512--1525] The original node and every neighbor are deep-copied.
Trial candidates temporarily mutate real data, so restore must recover arrays
as well as scalar node fields.
\item[Lines 1528--1536] Original \(f\mathbin{\cdot}v\) power is the baseline.
A split is attractive only if post-split motion releases more power.
\item[Lines 1539--1568] \code{SMNGetArmSets()} enumerates unique nontrivial
partitions. Position/velocity arrays remember the best trial, while
\code{armList} is built in descending arm-index order.
\item[Lines 1569--1609] Each partition is applied with \code{globalOp=0};
the zero-length trial link is local and temporary. Selected arms attach to
\code{splitNode2}. For BCC degree-three work, the binary-junction arm cannot
be put in the selected subset.
\item[Lines 1611--1642] Existing segment forces are copied from backups,
avoiding a full force calculation at the zero-length trial. The two temporary
nodes then receive mobility evaluation, with inclusion support when enabled.
\end{description}

\lstinputlisting[style=cpp,firstline=1644,lastline=2001,firstnumber=1644]{../src/Topology.cc}

\paragraph{Line-by-line explanation.}
\begin{description}
\item[Lines 1644--1665] A successful mobility evaluation is adjusted for new
junction force/velocity. The initially zero-length connecting segment becomes
important once nodes separate.
\item[Lines 1667--1687] The BCC degree-three case freezes its resulting
degree-two corner during trial scoring. Squared speeds identify which node or
nodes should move.
\item[Lines 1688--1791] If post-split velocities separate endpoints, the
faster side determines a trial direction. Equal-and-opposite motion moves
both nodes half the initial distance; otherwise one moves the full distance.
\code{FoldBox()} maintains periodic coordinates.
\item[Lines 1800--1888] Forces are recalculated after geometry changes. The
new nonzero connector receives a valid glide plane before the second mobility
call. Relative velocity is then saved for the power test.
\item[Lines 1905--1920] Mobility failure terminates candidate evaluation for
this node. A candidate is physically separating only when relative velocity
has a positive component along selected separation direction.
\item[Lines 1922--2001] Trial power is the sum of post-split
\(f\mathbin{\cdot}v\) terms minus force-scaled noise. Only a strict
improvement becomes best. Its segment forces/positions are saved; a four-arm
split that would immediately collide is marked for rejection.
\end{description}

\lstinputlisting[style=cpp,firstline=2003,lastline=2396,firstnumber=2003]{../src/Topology.cc}

\paragraph{Line-by-line explanation.}
\begin{description}
\item[Lines 2003--2067] Every temporary split is merged back locally and the
original node/neighbors are restored from deep backups. This is a true
rollback, so original arm indices remain valid across trials.
\item[Lines 2079--2097] A candidate is discarded if predicted immediate
collision would undo the new configuration.
\item[Lines 2099--2217] The winning partition is rebuilt in descending order.
Its saved direction determines final positions. The code either uses uniform
separation or calls \code{CalcJunctionLen()} for a line-tension-based length.
\item[Lines 2219--2269] Final global \code{SplitNode()} creates graph
mutation. A degree-two node from the BCC three-arm path receives
\code{CORNER_NODE}, and its constraint reset is appended to the MPI log.
\item[Lines 2271--2345] Plastic strain is corrected when enabled. Saved
forces are restored to the winning configuration and pushed to every
affected endpoint with \code{ResetSegForces2()}.
\item[Lines 2349--2396] The two participants become collision-exempt for the
rest of the cycle. Candidate and backup arrays are then freed.
\end{description}

\subsection{BCC binary-junction classifier}

\lstinputlisting[style=cpp,firstline=1231,lastline=1333,firstnumber=1231]{../src/Topology.cc}

\paragraph{Line-by-line explanation.}
\begin{description}
\item[Lines 1231--1264] The classifier returns an arm index or \(-1\). It
requires exactly three arms, initializes a candidate junction direction, and
uses a five-degree geometric tolerance.
\item[Lines 1266--1308] Each nonzero-length arm direction is calculated in
the nearest periodic image. An arm with Burgers magnitude above one is treated
as the junction arm; the other two supply edge/glide-plane normals.
\item[Lines 1311--1329] A valid binary junction has one junction arm and two
glissile arms. Cross products of parent normals decide whether the junction
line is planar. Planar cases are excluded by the caller; the extension splits
the non-planar case and preserves its physical corner.
\item[Lines 1332--1333] The junction arm index lets the caller reject
partitions that would move it to the degree-two side.
\end{description}

\subsection{Parallel multi-node splitting}

\lstinputlisting[style=cpp,firstline=15,lastline=123,firstnumber=15]{../src/SplitMultiNodesParallel.cc}

\paragraph{Line-by-line explanation.}
\begin{description}
\item[Lines 17--34] The same cadence policy applies before allocating
parallel-work data.
\item[Lines 36--53] The domain builds ghost and native multi-node lists. It
evaluates partial forces for every candidate partition because remote
contributions are needed to score local candidates accurately.
\item[Lines 55--71] Ghost-candidate force data is packed for owning domains.
The code determines which domains should send contributions for native
candidates, then exchanges buffers.
\item[Lines 73--99] All transient communication buffers are freed before the
final decision stage because their contents have been accumulated into native
multi-node records.
\item[Lines 101--106] \code{SMNEvalMultiNodeSplits()} now has complete forces
for local candidates and commits only locally owned winners. The later generic
operation exchange propagates resulting graph changes.
\item[Lines 108--123] Candidate storage is released and the shared timer
stops.
\end{description}

The records in \srcfile{include/SMN.h}, lines 43--71 retain the number of
arms and per-arm endpoint forces for each trial. Parallel mode changes where
forces are evaluated, not the physical acceptance criterion or the
\code{SplitNode()} graph operation.

\subsection{Cross slip and collision dispatch}

\lstinputlisting[style=cpp,firstline=636,lastline=656,firstnumber=636]{../src/CrossSlip.cc}

\lstinputlisting[style=cpp,firstline=942,lastline=954,firstnumber=942]{../src/Collision.cc}

\paragraph{Line-by-line explanation.}
\begin{description}
\item[\srcfile{CrossSlip.cc}, lines 636--640] The dispatcher reads run-time
enable state, material type, and BCC model selector.
\item[\srcfile{CrossSlip.cc}, lines 642--654] When enabled, exactly one
material-specific routine runs. These routines can reset glide planes and
reposition nodes; they use the current topology operation list.
\item[\srcfile{Collision.cc}, lines 942--951] \code{collisionMethod} selects
predictive-LG (4), retroactive (3), predictive (2), or proximity (1/default)
search/response logic. The default parameter value is 2.
\item[\srcfile{Collision.cc}, lines 953--954] The dispatcher has no geometry;
it isolates policy selection from callers.
\end{description}

\lstinputlisting[style=cpp,firstline=776,lastline=1040,firstnumber=776]{../src/PredictiveCollision.cc}

\paragraph{Line-by-line explanation of predictive-collision selection.}
\begin{description}
\item[Lines 776--828] The implementation declares four endpoints, positions,
velocities, candidate split/merge records, and \code{rann^2} threshold.
\item[Lines 830--887] Eligible native nodes visit their own and neighboring
\code{cell2} queues. \code{NO_COLLISIONS}, absent membership, and stable
\code{CollisionNodeOrder()} eliminate invalid or duplicate work.
\item[Lines 889--958] Nested arm loops form segment pairs. Current domain
must own both segments before it can change endpoints. Shared-hinge and
remote-owned candidates are skipped.
\item[Lines 960--1004] Four endpoints are mapped into consistent periodic
images and zero-length segments are rejected before distance mathematics.
\item[Lines 1006--1028] \code{GetMinDist()} calculates closest approach,
approach rate, and fractional positions. Intersections closer than one
Burgers-vector unit are treated as current collisions.
\item[Lines 1029--1040] Nearby nonintersecting segments enter predictive
time/point calculation. The remainder refines affected segments when needed,
then calls \code{MergeNode()} at a permitted point, updates forces, and marks
participants exempt. It uses the same split/merge primitives, not a separate
graph-edit system.
\end{description}

\section{MPI topology synchronization}

\subsection{Operation-log protocol}

\lstinputlisting[style=cpp,firstline=14,lastline=103,firstnumber=14]{../include/OpList.h}

\lstinputlisting[style=cpp,firstline=45,lastline=114,firstnumber=45]{../src/RemoteOps.cc}

\lstinputlisting[style=cpp,firstline=670,lastline=708,firstnumber=670]{../src/RemoteOps.cc}

\paragraph{Line-by-line explanation.}
\begin{description}
\item[\srcfile{OpList.h}, lines 14--28] \code{OpType_t} enumerates every
serialized mutation: connection changes, arm insertion/deletion through
Burgers change, node split/removal, coordinate/force/plane/constraint reset.
All ranks must run a binary with matching record layouts.
\item[\srcfile{OpList.h}, lines 30--103] Each record begins with
\code{opType}, followed by tags and scalar/vector payload to replay exactly
one mutation. Tags, rather than pointers, make a byte stream meaningful on a
different rank.
\item[\srcfile{RemoteOps.cc}, lines 45--51] Clearing an operation list zeroes
used bytes and resets used length. It is called once before physics topology
work and again before remesh.
\item[\srcfile{RemoteOps.cc}, lines 90--114] The list begins as a 4096-byte
buffer and grows in fixed increments. Producers need not predict count.
\item[\srcfile{RemoteOps.cc}, lines 670--708] A split producer enters an
OpenMP critical section, ensures capacity, fills the fixed record (old tag,
new tag, coordinates, velocity), copies it to the byte buffer, and advances
\code{opListUsedLen}. Other \code{AddOp*()} helpers use this append pattern.
\end{description}

\subsection{Sending the log}

\lstinputlisting[style=cpp,firstline=21,lastline=136,firstnumber=21]{../src/CommSendRemesh.cc}

\paragraph{Line-by-line explanation.}
\begin{description}
\item[Lines 23--32] The routine finds remote neighbor count and current byte
length of the local operation log.
\item[Lines 33--43] In a parallel build, nonblocking receives for every
neighbor's message length are posted first. This avoids send/send deadlock.
\item[Lines 45--70] MPI cannot safely use one buffer for multiple
nonblocking sends. The code retains the original for one receiver and clones
it for every other remote domain.
\item[Lines 72--100] Message lengths are exchanged and incoming byte buffers
are allocated before full record streams are received.
\item[Lines 102--117] Actual streams are sent/received asynchronously; all
requests complete before return.
\item[Lines 119--135] Temporary outgoing clones are released. The primary
\code{home->opListBuf} remains owned by \code{Home_t}; serial builds make the
network part a no-op.
\end{description}

\subsection{Replaying neighbouring operations}

\lstinputlisting[style=cpp,firstline=25,lastline=157,firstnumber=25]{../src/FixRemesh.cc}

\lstinputlisting[style=cpp,firstline=158,lastline=313,firstnumber=158]{../src/FixRemesh.cc}

\paragraph{Line-by-line explanation.}
\begin{description}
\item[Lines 25--55] For every remote domain, \code{FixRemesh()} reads its
received buffer from offset zero until the advertised length is consumed. The
first integer chooses a record format.
\item[Lines 57--90] Connection-change and insert-arm records are copied from
the stream, the offset advances by \code{sizeof(record)}, and corresponding
local endpoints are changed with \code{globalOp=0}. Replay must not append a
new operation recursively.
\item[Lines 92--157] A split record creates a ghost node at the specified
tag, coordinates, and velocity, frees stale arms, and assigns a spatial cell.
Later records add arms/connections. Producer ordering therefore matters:
split precedes records that refer to its new tag.
\item[Lines 158--212] Force-obsolete and force-reset records preserve force
consistency. They are not cosmetic; later geometric selection may refuse a
segment when forces are stale.
\item[Lines 214--247] Coordinate and constraint records change persistent
state. Constraint replay is how \code{CORNER_NODE} survives domain boundaries.
\item[Lines 249--299] Node removal, Burgers/arm changes, and glide-plane reset
records invoke the same low-level helpers used for local work, with logging
disabled.
\item[Lines 301--313] An unknown opcode is fatal because safe byte-offset
recovery is impossible. Once a domain's stream is replayed its buffer is
released.
\end{description}

\section{Remeshing}

\subsection{Dispatcher and postconditions}

\lstinputlisting[style=cpp,firstline=1172,lastline=1272,firstnumber=1172]{../src/Remesh.cc}

\paragraph{Line-by-line explanation.}
\begin{description}
\item[Lines 1174--1181] A negative \code{remeshRule} disables this pass.
Otherwise the dispatcher obtains parameters and may synchronize timer
measurement at the start.
\item[Lines 1190--1204] It starts the remesh timer, clears a fresh operation
log, resets per-cycle exemptions, and dispatches rule 2 or rule 3. Any other
nonnegative rule is fatal rather than silently selecting an algorithm.
\item[Lines 1207--1217] Free-surface builds remove a segment that joins two
nodes on the same surface.
\item[Lines 1219--1231] The second operation-log exchange propagates all
remesh splits/merges and replays remote changes.
\item[Lines 1233--1245] Independent remote changes may create a duplicate
link. Every local node is checked after replay.
\item[Lines 1247--1252] Orphaned nodes are removed before the next ghost-node
exchange. This is the final graph-cleanliness postcondition.
\item[Lines 1254--1271] Optional timing/memory checks close the pass.
\end{description}

\subsection{Force estimates used during topology changes}

\lstinputlisting[style=cpp,firstline=132,lastline=167,firstnumber=132]{../src/Remesh.cc}

\lstinputlisting[style=cpp,firstline=187,lastline=290,firstnumber=187]{../src/Remesh.cc}

\paragraph{Line-by-line explanation.}
\begin{description}
\item[Lines 132--167] \code{EstRefinementForces()} finds the existing arm
indices, copies Burgers data and endpoint arm forces, represents node 2 in the
same periodic image as node 1, then calls \code{FindSubFSeg()} to estimate
four endpoint forces on the two child segments.
\item[Lines 187--224] \code{EstCoarsenForces()} finds both old arm IDs and
makes all three nodes periodic-image consistent. It copies Burgers and force
data needed for a combined future segment.
\item[Lines 225--290] Very short/zero arms are treated specially because
\code{FindFSegComb()} is not stable for them. Otherwise it estimates endpoint
forces for the one segment that replaces two old segments.
\end{description}

These are estimates, not a substitute for full force calculation. Later
cleanup calls \code{MarkNodeForceObsolete()} so accurate forces are calculated
at the next appropriate force-update point.

\subsection{Mesh coarsening: selecting a removal}

\lstinputlisting[style=cpp,firstline=378,lastline=419,firstnumber=378]{../src/Remesh.cc}

\lstinputlisting[style=cpp,firstline=460,lastline=758,firstnumber=460]{../src/Remesh.cc}

\paragraph{Line-by-line explanation of candidate selection.}
\begin{description}
\item[Lines 378--419] \code{MeshCoarsen()} derives local/remote maximum-length
cutoffs, squared area threshold, and allocates candidate storage. The list
allows the initial geometry scan to be parallel while graph changes are
serialized safely.
\item[Lines 460--495] A candidate must be degree two, unpinned, not already
exempt, have two resolvable neighbors, and own at least one incident segment.
\item[Lines 497--528] The two arms and future neighbor-to-neighbor chord are
computed in minimum-image coordinates as \(r_1,r_2,r_3\).
\item[Lines 530--543] \code{CORNER_NODE} protects a physical junction corner.
It can be removed only across a sufficiently short attached segment, avoiding
loss of a real kink merely because it is degree two.
\item[Lines 545--651] When glide planes are enforced, the arms normally must
share a plane. Exceptions allow debris cleanup for tiny segments, small
triangles, fuzzy/close planes, and matching nearest precise planes.
\item[Lines 653--668] The future direct segment cannot exceed \code{maxSeg}.
With a remote neighbor, the stricter half-cell cutoff prevents two independent
domain changes creating a segment spanning too many cells.
\item[Lines 671--715] Heron's formula gives \code{area2}, which is actually
\(A^2\), not twice the area. Its time derivative determines shrinking. A
candidate is listed when \(A^2<A_{\min}^2\) and shrinking, or either arm is
shorter than \code{minSeg}.
\item[Lines 717--758] Unsafe remote-ownership cases are rejected. Accepted
candidates are appended under a critical section as \code{REMESH_PENDING}.
\end{description}

\subsection{Mesh coarsening: committing nonoverlapping removals}

\lstinputlisting[style=cpp,firstline=764,lastline=1051,firstnumber=764]{../src/Remesh.cc}

\lstinputlisting[style=cpp,firstline=1053,lastline=1160,firstnumber=1053]{../src/Remesh.cc}

\paragraph{Line-by-line explanation of commit and repair.}
\begin{description}
\item[Lines 764--816] Candidates are revisited in parallel, but
\code{ProceedWithCoarsen()} runs in a critical section. It prevents two
events sharing a node or neighbor from both committing.
\item[Lines 818--838] Current nodes are resolved from saved tags because
earlier events can delete one. Missing data safely cancels the event.
\item[Lines 840--885] Corner-protected endpoints may only merge over a link
no longer than \(2rann\). The code tests both potential target neighbors.
\item[Lines 887--918] Segment forces are estimated before mutation. An event
is global if either neighbor, or any arm connected to either neighbor, reaches
another domain.
\item[Lines 926--983] Actual mutation is in a critical section. Code first
attempts \code{MergeNode(node,nbr1)} at \code{nbr1}'s location, then tries
the other neighbor if the first merge is illegal.
\item[Lines 985--1051] Surviving endpoints are resolved again, estimated
forces are installed on the new segment, and the result becomes
coarsen-exempt during this pass.
\item[Lines 1053--1139] A stable-topology pass marks endpoint and neighbor
forces obsolete. When fuzzy planes are enabled, new segment plane is
recalculated.
\item[Lines 1141--1160] Optional diagnostics reduce counts across ranks,
then temporary candidate storage is freed.
\end{description}

\subsection{Rule 2 refinement: midpoint insertion}

\lstinputlisting[style=cpp,firstline=35,lastline=65,firstnumber=35]{../src/RemeshRule_2.cc}

\lstinputlisting[style=cpp,firstline=95,lastline=420,firstnumber=95]{../src/RemeshRule_2.cc}

\paragraph{Line-by-line explanation of selection.}
\begin{description}
\item[Lines 35--65] The function derives squared maximum length/area
thresholds and allocates a refine list. The leading source comment calls this
version 3, but dispatcher lines in \srcfile{Remesh.cc} show that the file
implements rule 2; this is a source-comment typo.
\item[Lines 95--104] Pinned nodes are never used as degree-two refinement
anchors.
\item[Lines 106--179] A degree-two node uses its two arm lengths and local
triangle geometry. Rule 2 area derivative includes all endpoint velocities.
\item[Lines 181--217] Each arm starts eligible, then is rejected if forces at
both endpoints are stale or this domain does not own the segment.
\item[Lines 219--318] A segment is selected if it exceeds \code{maxSeg}, or
if it is sufficiently curved: it is not too short to bisect, local area
exceeds \code{remeshAreaMax}, and local area is not shrinking. Duplicate
checks ensure only one list entry represents the physical segment.
\item[Lines 320--420] Nodes whose degree is not two use a simpler rule: the
stable lower tag endpoint considers each owned arm once and lists it only when
its squared length exceeds \code{maxSeg^2}.
\end{description}

\lstinputlisting[style=cpp,firstline=426,lastline=661,firstnumber=426]{../src/RemeshRule_2.cc}

\paragraph{Line-by-line explanation of Rule 2 commit.}
\begin{description}
\item[Lines 426--479] Saved endpoint tags are resolved and both nodes lock.
Stored segment vector puts the point at the midpoint; endpoint velocities are
averaged; replacement-segment forces are estimated before arm arrays change.
\item[Lines 481--526] Insertion point is folded into the periodic box. A split
is global only across domains. Under a critical section,
\code{SplitNode()} moves exactly one arm from node 1 to a new node, replacing
\(n_1--n_2\) with \(n_1--n_{new}--n_2\).
\item[Lines 528--548] Locks are released. A successful entry saves new tag
and four estimated endpoint-force vectors.
\item[Lines 551--618] Each committed event marks three nodes force-stale,
installs estimates on the two new segments, and reevaluates mobility under
the relevant node lock.
\item[Lines 642--661] Optional logging prints transformation and list storage
is released.
\end{description}

\subsection{Rule 3 refinement: fitted-arc insertion}

\lstinputlisting[style=cpp,firstline=123,lastline=377,firstnumber=123]{../src/RemeshRule_3.cc}

\paragraph{Line-by-line explanation of Rule 3 geometry.}
\begin{description}
\item[Lines 123--190] Rule 3 starts with the same degree-two triangle, but
computes area rate using only the central node velocity. Treating neighbors as
stationary prevents choosing an arc direction merely because neighbors move
away while the central node moves toward them.
\item[Lines 192--231] Ownership and stale-force guards are unchanged from
rule 2.
\item[Lines 233--267] Two segment midplanes and normal
\(\vect{v}_1\times\vect{v}_2\) define three equations for the circumcentre of
the local three-node arc.
\item[Lines 269--313] Rule 3 uses the same basic length/area decision as rule
2 before spending work on fitted geometry.
\item[Lines 315--323] For sufficiently large area, the \(3\times3\) matrix is
inverted to calculate the circle centre. An inversion failure is fatal; the
source intentionally does not fall back to a midpoint in this branch.
\item[Lines 325--339] The radius and radial direction through the selected
chord midpoint are calculated.
\item[Lines 341--356] The segment glide-plane normal is calculated as
\(\vect{b}\times\vect{l}\). A screw segment has zero normal, which is handled
without division by zero.
\item[Lines 358--377] Midpoint is moved onto the circumcircle, then its
normal-to-glide-plane displacement is removed. The final \code{newPos}
preserves local curvature while respecting the plane. Below curvature
threshold, ordinary midpoint insertion is retained.
\end{description}

\lstinputlisting[style=cpp,firstline=419,lastline=434,firstnumber=419]{../src/RemeshRule_3.cc}

\lstinputlisting[style=cpp,firstline=560,lastline=785,firstnumber=560]{../src/RemeshRule_3.cc}

\paragraph{Line-by-line explanation of Rule 3 commit.}
\begin{description}
\item[Lines 419--434] The computed arc position is stored with
\code{haveNewPos=1}. The record also retains the original segment vector for
force estimation and fallback processing.
\item[Lines 560--606] At commit, \code{haveNewPos} selects stored arc point or
midpoint. Velocity interpolation and force estimation remain as in rule 2.
\item[Lines 608--673] Global-operation decision, one-arm \code{SplitNode()}
call, locking, and committed-record update match Rule 2 semantics.
\item[Lines 676--785] Final force-obsolete, estimated-force, and mobility
updates also match. The rules differ primarily in candidate geometry and
insertion position, not graph mutation or MPI behavior.
\end{description}

\subsection{Rule comparison}

\begin{center}
\begin{tabularx}{0.96\linewidth}{@{}lXX@{}}
\toprule
Rule & Selection trigger & Insertion placement \\
\midrule
2 & Long segment or a nonshrinking high-area degree-two triangle & Midpoint of selected segment. \\
3 & Same basic trigger, with extra degree-two curved-geometry work & Fitted circumcircle/arc point projected to glide plane; midpoint fallback. \\
\bottomrule
\end{tabularx}
\end{center}

Both rules call \code{MeshCoarsen()} first, then refine. Both use the same
operation records, locks, force estimates, MPI replay, duplicate-link cleanup,
and orphan removal.

\section{End-to-end transactions}

\subsection{Example: a long degree-two segment}

\begin{enumerate}
\item \code{ParadisStep()} reaches \code{Remesh()} only after the
physics-driven topology exchange completes.
\item \code{Remesh()} clears operation log, resets exemptions, and calls
\code{RemeshRule_2()} or \code{RemeshRule_3()}.
\item \code{MeshRefine()} recognizes an owned segment whose length exceeds
\code{maxSeg}. Rule 2 chooses a midpoint; rule 3 may choose an arc point for
a curved degree-two configuration.
\item It locks endpoints, estimates forces for two replacements, then calls
\code{SplitNode(OPCLASS_REMESH,...)} with one selected arm.
\item \code{SplitNode()} creates a native node, relinks old neighbor, inserts
the required connecting arm, and appends operations if the segment crosses
domains.
\item Refinement installs estimates and marks nodes for accurate future force
calculation. \code{CommSendRemesh()} and \code{FixRemesh()} make remote ghost
topology match; cleanup removes duplicates and orphans.
\end{enumerate}

\subsection{Example: a short, collapsing degree-two node}

\begin{enumerate}
\item \code{MeshCoarsen()} verifies degree two, unpinned status, ownership,
glide-plane compatibility, maximum future length, and corner protections.
\item A shrinking small triangle or an arm shorter than \code{minSeg} adds a
candidate. Overlap control chooses a nonoverlapping subset.
\item The code estimates future direct-segment force and calls
\code{MergeNode(OPCLASS_REMESH,...)} at an allowed neighbor position.
\item \code{MergeNode()} chooses legally deletable endpoint, transfers arms
to survivor, inherits constraints, and invokes duplicate-link cleanup.
\item The new segment receives estimated forces; exact forces are marked
obsolete. The second MPI exchange propagates deletion and relinking.
\end{enumerate}

\subsection{Example: a multi-node separation}

\begin{enumerate}
\item The serial path enumerates arm partitions and makes each a temporary
zero-length \code{SplitNode()} transaction. Parallel path calculates the same
candidate force data across contributing domains.
\item Mobility, junction force, separation direction, and
\(f\mathbin{\cdot}v\) power select the best admissible partition.
\item Every trial is rolled back. Only winning split runs with
\code{globalOp=1}, writes operation records, installs saved forces, and marks
participants collision-exempt.
\item Collision handling then runs and replays full log remotely. A later
remesh pass can refine a new junction segment but does not normally coarsen a
protected BCC physical corner.
\end{enumerate}

\section{Validation, limitations, and developer checklist}

\subsection{Initialization checks that protect remesh}

\lstinputlisting[style=cpp,firstline=690,lastline=739,firstnumber=690]{../src/Initialize.cc}

\paragraph{Line-by-line explanation.}
\begin{description}
\item[Lines 692--702] \code{minSeg} must be less than half \code{maxSeg}, and
\code{maxSeg} must exceed \code{minSeg}. This assures that one bisection can
produce acceptable child lengths.
\item[Lines 704--709] \code{maxSeg} must remain under 90 percent of the
smallest cell width, preserving force/cell spatial assumptions.
\item[Lines 711--715] Derived area thresholds must remain mutually
consistent; bad geometry tuning is stopped before stepping.
\item[Lines 717--739] Nonfatal warnings expose poor core radius, area, and
contact-distance choices that can destabilize topology decisions.
\end{description}

\subsection{Implementation details to keep in mind}

\begin{itemize}
\item In \srcfile{RemeshRule_2.cc}:227 and
\srcfile{RemeshRule_3.cc}:274, \code{segLen} is declared \code{int} but is
assigned a real-valued arm length before comparisons. Consequently the
degree-two length comparisons truncate fractional length. This guide records
current behavior; maintainers may wish to audit whether \code{real8} was
intended.
\item Candidate arrays are initially sized at twice
\code{home->newNodeKeyPtr}. If full, append code skips later candidates rather
than reallocating. That is safe but can defer legitimate remesh work.
\item Rule 3 treats a circumcentre-matrix inversion failure as fatal. This
avoids silently inserting at invalid position, but it makes nearly degenerate
curved cases worth testing.
\item Cross slip is in topology phase because it changes segment planes and
positions, but it normally preserves node/segment cardinality. Do not describe
all cross-slip events as node creation or deletion.
\end{itemize}

\subsection{Invariants for topology modifications}

\begin{itemize}
\item \textbf{Ownership:} mutate a segment only when \code{DomainOwnsSeg()}
grants operation-class ownership; mutate/delete a node only through legal
primitive checks.
\item \textbf{Burgers conservation:} \code{SplitNode()} creates a connector
from moving-side Burgers sum. Outer driver calls
\code{CheckBurgConservation()} after topology replay.
\item \textbf{Both endpoints agree:} never edit one node's \code{nbrTag},
Burgers, plane, or force array by hand. Use connection/arm helpers.
\item \textbf{No overlaps:} candidate detection may be parallel; commits must
lock/serialize shared nodes and maintain
\code{REMESH_COMMITTED}/\code{REMESH_REJECTED}.
\item \textbf{Force validity:} respect \code{NODE_RESET_FORCES}, preserve
estimates, and mark affected neighborhoods obsolete after mutation.
\item \textbf{Periodic geometry:} calculate vectors with \code{ZImage()} or
\code{PBCPOSITION()}, and fold stored coordinates with \code{FoldBox()}.
\item \textbf{Distributed replay:} every remote-visible change needs an
\code{AddOp*()} record. Replay it with \code{globalOp=0}.
\item \textbf{Postconditions:} clean duplicate links after remote replay and
remove local orphan nodes before ghost communication.
\end{itemize}

\subsection{Useful breakpoints and logging}

{\small
\begin{longtable}{@{}p{0.38\linewidth}p{0.56\linewidth}@{}}
\toprule
Location or symbol & What it reveals \\
\midrule
\endhead
\code{src/ParadisStep.cc:403,419,447,493} & High-level phase order for one cycle. \\
\code{src/Topology.cc:2465} \code{SplitNode()} & Every arm-moving insertion, including remesh. \\
\code{src/Topology.cc:2976} \code{MergeNode()} & Every coarsen/collision deletion decision. \\
\code{src/Remesh.cc:714} and \code{src/RemeshRule_2.cc:254} & Geometry predicates that nominate coarsen/refine candidates. \\
\code{src/RemoteOps.cc:670} and \code{src/FixRemesh.cc:92} & A split record being produced and replayed. \\
\code{DEBUG_TOPOLOGY_CHANGES} & Prints separation, coarsen, refine, and collision operations. \\
\code{DEBUG_LOG_MESH_COARSEN}, \code{DEBUG_LOG_MESH_REFINE} & Reduces and prints event counts across ranks. \\
\end{longtable}
}

\subsection{Checklist for a new topology event}

\begin{enumerate}
\item Decide whether event moves arms, merges nodes, changes a segment
property, or needs a new primitive. Reuse \code{SplitNode()} and
\code{MergeNode()} whenever semantics fit.
\item Define deterministic ownership, including a safe response when a
candidate spans domains.
\item Screen candidates using force validity, constraints, periodic geometry,
and per-cycle exemptions.
\item Make trial operations local/reversible; make final operation global only
when remote state can observe it.
\item Record every remote-visible change through an \code{AddOp*()} operation,
or add record type, producer, and \code{FixRemesh()} case together.
\item Restore/recompute per-arm forces and mobility; mark untrusted data
obsolete.
\item Exercise serial, OpenMP, and MPI cases. Verify Burgers conservation,
duplicate-link cleanup, and restart/ghost consistency.
\end{enumerate}

{\footnotesize\noindent\textit{Primary sources: the Section~1 code files;
\code{docs/ParaDiS_4.0_UserGuide.pdf}; and
\code{docs/BCC_Binary_Junction_Node_Splitting_code.tex}.}\par}

\end{document}
